\item La “ola” es un tipo particular de pulso que se puede propagar a través de una gran multitud reunida en un estadio deportivo (figura P16.4). Los elementos del medio son los espectadores, con posición cero cuando están sentados y posición máxima cuando están de pie y elevan sus brazos. Si una gran cantidad de espectadores participa en el movimiento ondulatorio, se desarrolla una forma de pulso estable. La rapidez de la onda depende del tiempo de reacción de las personas, que por lo general es del orden de 0.1 s. Estime el orden de magnitud, en minutos, del intervalo de tiempo requerido para que tal pulso dé una vuelta completa en torno a un gran estadio deportivo. Establezca las cantidades que mida o estime y sus valores.